\section{Experimental Setup}
\label{sec:setup}

\paragraph{Models.}
Our main model is Qwen3-VL-8B-Instruct~\cite{qwen3vl}. It sees 64 frames at native resolution, sampled uniformly
from $[0,q]$ or chosen by BOLT (\cref{sec:retrieval}), and decodes greedily or with PCD, at batch size 8 for pruned
memories. We generate at most 384 new tokens on S-EMBER (repetition penalty 1.1), 512 on TimeLens-Bench, and 256
(two-query) or 1{,}024 (joint answer) on NExT-GQA, both with repetition penalty 1.0. All methods share the timeline
prompt, except on the full TimeLens-Bench sets (including the QVHighlights and Charades-STA columns of
\cref{tab:main}) and in the prompt study of \cref{sec:transfer}, which use the TimeLens-Bench one-line prompt.
\cref{sec:transfer} also adds TimeLens-8B~\cite{timelens}, a Qwen3-VL-8B model fine-tuned for grounding, with
frames capped at about 1.57\,MP on S-EMBER and NExT-GQA. On TimeLens-Bench both models use the protocol input
(2\,fps, at most 14{,}336 visual tokens and at least 64 per frame), and BOLT runs use 64 or 128 frames of at most
0.79\,MP. Batch size alone changes about half of the greedy outputs and moves a score by up to about two points, so
comparisons with the uncompressed model (batch size 1--2) include this noise.

\paragraph{Two clusters.}
Later runs used a second GPU cluster: most fixed-rule and all 5\% NExT-GQA runs, BOLT on the full QVHighlights
set, the official S-EMBER prompt, and the TimeLens-8B gold-informed runs. Paired differences that combine both
clusters are marked. Runs repeated on both clusters agree within one mIoU point, \eg $+0.4$ for NExT-GQA without
pruning (469 of 510 replies identical) and $-0.9$ for \method on S-EMBER.

\paragraph{Datasets.}
\textbf{S-EMBER} grounded QA~\cite{sember}: 475 questions (191 duration, 184 counting, 100 location) on the first
300 videos, asked mid-stream, so the memory holds only the frames before the question.
\textbf{QVHighlights}~\cite{qvhighlights} val: 500 queries with a single relevant window, sampled with seed 0 from
the 1{,}020 single-window queries, on videos of about 150\,s.
\textbf{ActivityNet Captions}~\cite{activitynet} val\_1: 500 queries, one sentence per video, on videos of about
115\,s. \textbf{Charades-STA}~\cite{charadessta} test: 500 queries on videos of about 30\,s. For these three datasets
the question is \emph{``When does the following happen in the video: \textless query\textgreater?''} and the window is
the whole video.
\textbf{TimeLens-Bench}~\cite{timelens}: the re-annotated Charades-STA, QVHighlights and ActivityNet sets, with their
one-line prompt, on the full sets and on a 513-query QVHighlights subset. \textbf{NExT-GQA}~\cite{nextgqa} test:
1{,}005 questions on videos of about 40\,s, with the answer and the interval asked in two separate queries (multiple
choice or open-ended) or, for open-ended answers, in one joint reply. Our QVHighlights and ActivityNet copies are
low-resolution re-encodes with 153 and about 84 visual tokens per frame pair, \vs 660 on S-EMBER, so we also report
budgets in absolute tokens.

\paragraph{Budget.}
The default budget is 10\% of the visual tokens of the 64-frame memory, timestamps included (about 2{,}140 tokens on
S-EMBER, about 520 on QVHighlights). ActivityNet is reported at 30\% (400--900 tokens), because at 10\% its memory
holds only 120--300 tokens (\cref{sec:budget}). Since timestamps count against the budget, a KV selector keeps fewer
visual tokens than an input pruner at the same nominal $\rho$, about 3.4\% on NExT-GQA at 10\%. A frame's timestamp
is kept when the frame keeps a visual token. An earlier version of this rule could let timestamps fill the budget of
very small memories and leave no visual token. The NExT-GQA 10\% runs use the fixed rule (\cref{sec:transfer}) and
the other runs the earlier one, which can matter only where timestamps take a large share of the budget
(\cref{sec:budget}).

\paragraph{Baselines.}
We compare against all tokens (no pruning), random tokens, and four QA pruners: FastVID~\cite{fastvid},
VidCom\textsuperscript{2}~\cite{vidcom2}, FlashVID~\cite{flashvid} (vision-side stages) and SemVID~\cite{semvid},
re-implemented and checked to keep the same tokens as the authors' code on test inputs. Each baseline is run with
and without PCD. A global question top-$k$ without quotas is compared in \cref{sec:res-coverage}.
\TODO{SnapKV-style observation-window eviction, PruneVid, and FV-Action / attention readout on the same memory.}

\paragraph{Metrics.}
We report mIoU and R@$\{0.3,0.5,0.7\}$. On S-EMBER we also report answer accuracy and GQ@0.5, judged by a local
Qwen3.8-27B-FP8 model with the official S-EMBER judge prompt. Verdicts for the same run differ by up to about one accuracy point across our three judging sessions (\eg 21.7
in A \vs 21.9 in B for one \method run), so we never pair Acc.\ or GQ@0.5 across sessions. Each table uses one session, named
in its caption: A for \cref{tab:qa}, B for \cref{tab:tl8b,tab:oracle}, and C for \cref{tab:oracle-tl8b}. The
early-phase tables listed in \cref{sec:results} use their own earlier sessions and are compared only within
themselves. Open-ended NExT-GQA answers are judged by the same model with the Video-ChatGPT
prompt~\cite{videochatgpt}, in separate batches that share a cache. On S-EMBER, times given as minutes:seconds are
converted to seconds for every method (\cref{sec:suppl-parse}). The joint NExT-GQA answers are scored with the
official S-EMBER parser, which counts such times as missing.

\paragraph{Protocol.}
We made all design choices (scores, quotas, $\alpha$, scope) on S-EMBER, then froze the configuration and ran it
once on QVHighlights, ActivityNet and Charades-STA. \TODO{held-out S-EMBER videos as a test split.}
